% ============================================================
%  PEMBAHASAN SUPER DETAIL — LATIHAN SOAL & STUDI KASUS
%  BAB 9: BUNGA TUNGGAL DAN BUNGA MAJEMUK
%  Catatan: Soal Akhir Bab TIDAK dibahas di sini (sesuai permintaan).
% ============================================================
% ============================================================
%  PREAMBUL BERSAMA — BUKU PEMBAHASAN LATIHAN SOAL
%  Dipakai oleh MAE-2026-2027-Latihan1.tex ... Latihan9.tex
%  Kompilasi: pdfLaTeX (disarankan Overleaf / MiKTeX)
% ============================================================
\documentclass[11pt,a4paper]{report}

% ---------- PAKET ----------
\usepackage[utf8]{inputenc}
\usepackage[T1]{fontenc}
\usepackage[bahasa]{babel}
\usepackage{amsmath,amssymb,amsthm,mathtools}
\usepackage{geometry}
\usepackage{xcolor}
\usepackage{graphicx}
\usepackage{enumitem}
\usepackage{tikz}
\usetikzlibrary{calc,arrows.meta,angles,quotes,positioning,decorations.pathreplacing,patterns}
\usepackage{pgfplots}
\pgfplotsset{compat=1.18}
\usepackage[most]{tcolorbox}
\usepackage{fancyhdr}
\usepackage{titlesec}
\usepackage{array,booktabs}
\usepackage{multicol}
\usepackage{pifont}
\usepackage[colorlinks=true,linkcolor=biru,urlcolor=biru]{hyperref}

\geometry{a4paper,margin=2.5cm,headheight=15pt}

% ---------- WARNA ----------
\definecolor{biru}{HTML}{1F4E79}
\definecolor{birumuda}{HTML}{D6E4F0}
\definecolor{hijau}{HTML}{2E7D32}
\definecolor{hijaumuda}{HTML}{E3F2E1}
\definecolor{oranye}{HTML}{C75B12}
\definecolor{oranyemuda}{HTML}{FBE8DA}
\definecolor{abu}{HTML}{F2F2F2}
\definecolor{ungu}{HTML}{6A1B9A}
\definecolor{ungumuda}{HTML}{EDE2F3}
\definecolor{merah}{HTML}{C62828}
\definecolor{merahmuda}{HTML}{FCE4E4}
\definecolor{tosca}{HTML}{00838F}
\definecolor{toscamuda}{HTML}{E0F2F1}
\definecolor{emas}{HTML}{8D6E00}
\definecolor{emasmuda}{HTML}{FFF6D6}

% ---------- HEADER / FOOTER ----------
\newcommand{\babjudul}{Pembahasan Latihan}
\pagestyle{fancy}
\fancyhf{}
\renewcommand{\headrulewidth}{0.8pt}
\renewcommand{\footrulewidth}{0pt}
\fancyhead[L]{\small\textcolor{ungu}{\leftmark}}
\fancyhead[R]{\small\textcolor{ungu}{\babjudul}}
\fancyfoot[C]{\thepage}

% ---------- GAYA JUDUL BAB ----------
\titleformat{\chapter}[display]
  {\normalfont\huge\bfseries\color{ungu}}
  {\filright\large PEMBAHASAN}{8pt}{\Huge}
\titlespacing*{\chapter}{0pt}{0pt}{20pt}

% ---------- LINGKUNGAN TEOREMA ----------
\newtheoremstyle{gaya}{6pt}{6pt}{}{}{\bfseries\color{biru}}{.}{.5em}{}
\theoremstyle{gaya}
\newtheorem{definisi}{Definisi}[chapter]
\newtheorem{sifat}{Sifat}[chapter]

% ---------- KOTAK: SOAL (ungu) ----------
\newtcolorbox{soalbox}[1]{
  enhanced, breakable, colback=ungumuda, colframe=ungu, boxrule=0.8pt,
  arc=3pt, left=10pt, right=10pt, top=8pt, bottom=8pt,
  title={\textbf{Soal #1}}, fonttitle=\bfseries\color{white}, coltitle=white,
  attach boxed title to top left={xshift=10pt,yshift=-10pt},
  boxed title style={colback=ungu,arc=2pt}}

% ---------- KOTAK: KONSEP KUNCI (hijau) ----------
\newtcolorbox{ingat}[1]{
  enhanced, breakable, colback=hijaumuda, colframe=hijau, boxrule=0.6pt,
  arc=3pt, left=10pt, right=10pt, top=6pt, bottom=6pt,
  title={\textbf{Konsep Kunci: #1}}, fonttitle=\bfseries\color{white}, coltitle=white,
  attach boxed title to top left={xshift=10pt,yshift=-10pt},
  boxed title style={colback=hijau,arc=2pt}}

% ---------- KOTAK: JEBAKAN (merah) ----------
\newtcolorbox{jebakan}{
  enhanced, breakable, colback=merahmuda, colframe=merah, boxrule=0.6pt,
  arc=3pt, left=10pt, right=10pt, top=6pt, bottom=6pt,
  title={\textbf{\ding{55}\ Jebakan \& Kesalahan yang Sering Terjadi}},
  fonttitle=\bfseries\color{white}, coltitle=white,
  attach boxed title to top left={xshift=10pt,yshift=-10pt},
  boxed title style={colback=merah,arc=2pt}}

% ---------- KOTAK: VARIASI SOAL (oranye) ----------
\newtcolorbox{variasi}{
  enhanced, breakable, colback=oranyemuda, colframe=oranye, boxrule=0.6pt,
  arc=3pt, left=10pt, right=10pt, top=6pt, bottom=6pt,
  title={\textbf{\ding{72}\ Bentuk Modifikasi Soal Ini}},
  fonttitle=\bfseries\color{white}, coltitle=white,
  attach boxed title to top left={xshift=10pt,yshift=-10pt},
  boxed title style={colback=oranye,arc=2pt}}

% ---------- KOTAK: STUDI KASUS (biru) ----------
\newtcolorbox{studikasus}[1]{
  enhanced, breakable, colback=abu, colframe=biru, boxrule=0.8pt,
  arc=3pt, left=10pt, right=10pt, top=8pt, bottom=8pt,
  title={\textbf{Studi Kasus: #1}}, fonttitle=\bfseries\color{white}, coltitle=white,
  attach boxed title to top left={xshift=10pt,yshift=-10pt},
  boxed title style={colback=biru,arc=2pt}}

% ---------- KOTAK: TIPS (tosca) ----------
\newtcolorbox{tips}{
  enhanced, breakable, colback=toscamuda, colframe=tosca, boxrule=0.6pt,
  arc=3pt, left=10pt, right=10pt, top=6pt, bottom=6pt,
  borderline west={3pt}{0pt}{tosca}}

% ---------- PERINTAH BANTU ----------
\newcommand{\jawab}{\par\smallskip\noindent\textit{\textcolor{oranye}{Penyelesaian langkah demi langkah:}}\par\smallskip}
\newcommand{\langkah}[1]{\par\smallskip\noindent\textbf{\textcolor{biru}{Langkah #1.}}\ }
\newcommand{\jadi}{\par\smallskip\noindent$\Rightarrow$\ \textbf{Jadi, }}
% Judul tiap nomor soal
\newcommand{\soalno}[2]{%
  \par\bigskip
  \noindent{\large\bfseries\color{ungu}\rule[-2pt]{4pt}{14pt}\ Soal #1}%
  \ {\normalfont\itshape\color{gray}(#2)}\par\nopagebreak\smallskip}

% ---------- HALAMAN JUDUL OTOMATIS ----------
% #1 = judul topik (mis. "Barisan dan Deret"); #2 = nomor bab (angka)
\newcommand{\buathalamanjudul}[2]{%
  \begin{titlepage}
  \centering
  \vspace*{2cm}
  {\color{ungu}\rule{\textwidth}{2pt}}\\[0.6cm]
  {\Huge\bfseries\color{ungu} PEMBAHASAN\\[0.3cm] LATIHAN SOAL}\\[0.4cm]
  {\Large BAB #2 --- #1}\\[0.2cm]
  {\color{ungu}\rule{\textwidth}{2pt}}\\[1.2cm]

  {\large Matematika SMA --- Fase E (Kelas 10)}\\[0.2cm]
  {\large Kurikulum Merdeka}\\[2cm]

  \begin{tcolorbox}[enhanced,width=0.85\textwidth,colback=ungumuda,colframe=ungu,
    arc=4pt,boxrule=1pt]
  \centering\small
  Buku pembahasan ini mengupas \textbf{setiap soal pada kotak ungu (Latihan Soal)}
  di tiap subbab, \emph{selangkah demi selangkah seperti mengajari dari nol}:
  konsep yang dipakai, jebakan yang sering menjerat, cara soal bisa dimodifikasi,
  dan visualisasi. \emph{Soal Akhir Bab tidak dibahas di sini.}
  \end{tcolorbox}

  \vfill
  {\large Tahun Pelajaran 2026/2027}
  \end{titlepage}
  \tableofcontents
  \thispagestyle{empty}
  \clearpage}

% ============================================================
%  PENJAGA KOMPILASI
%  File ini adalah PREAMBUL BERSAMA (di-\input oleh Latihan1..9),
%  BUKAN untuk dikompilasi sendiri. Jika editor terlanjur mengompilasi
%  file ini secara langsung, tampilkan satu halaman info (bukan error).
%  Saat di-\input oleh berkas Latihan, blok ini dilewati.
% ============================================================
\ifnum\pdfstrcmp{\jobname}{MAE-2026-2027-Latihan-preamble}=0
  \begin{document}
  \thispagestyle{empty}
  \begin{center}
  \vspace*{3cm}
  {\Huge\bfseries\color{ungu} Berkas Preambel Bersama}\\[1cm]
  \begin{tcolorbox}[enhanced,width=0.85\textwidth,colback=ungumuda,
    colframe=ungu,arc=4pt,boxrule=1pt]
  \centering
  File ini \textbf{bukan untuk dikompilasi sendiri.} Ia berisi pengaturan
  bersama (paket, warna, kotak, perintah) yang dipanggil lewat
  \texttt{\textbackslash input} oleh berkas \textbf{MAE-2026-2027-Latihan1.tex}
  sampai \textbf{Latihan9.tex}.

  \medskip
  Silakan buka dan kompilasi salah satu berkas \textbf{Latihan1--Latihan9}
  tersebut untuk menghasilkan PDF pembahasan.
  \end{tcolorbox}
  \end{center}
  \end{document}
\fi

\renewcommand{\babjudul}{Pembahasan Latihan — Bab 9}

\begin{document}
\setcounter{chapter}{9}
\buathalamanjudul{Bunga Tunggal dan Bunga Majemuk}{9}

% ============================================================
\chapter*{Pembahasan Bab 9: Bunga Tunggal dan Bunga Majemuk}
\addcontentsline{toc}{chapter}{Pembahasan Bab 9: Bunga Tunggal \& Majemuk}
\markboth{PEMBAHASAN BAB 9}{}
% ============================================================

\begin{tips}
\textbf{Cara membaca buku ini.} Tiap soal: \textbf{(1)} soal ditulis ulang,
\textbf{(2)} konsep kunci, \textbf{(3)} langkah demi langkah, \textbf{(4)}
jebakan, \textbf{(5)} variasi. \emph{Aturan emas Bab 9:} tanyakan
``\textbf{tunggal} (bunga dari modal awal saja) atau \textbf{majemuk} (bunga
berbunga)?'', dan selalu \textbf{cocokkan satuan waktu} antara $i$ dan $n$.
\end{tips}

% ============================================================
\section*{Studi Kasus: Memilih Produk Tabungan}
\addcontentsline{toc}{section}{Studi Kasus: Memilih Produk Tabungan}
% ============================================================

\begin{studikasus}{Memilih Produk Tabungan}
Andi punya Rp$10.000.000$, menabung $3$ tahun. Bank A: bunga \textbf{tunggal}
$12\%$/tahun. Bank B: bunga \textbf{majemuk} $12\%$/tahun.
\begin{center}
\begin{tabular}{c|r|r}
\toprule
Tahun ke- & Bank A (Tunggal) & Bank B (Majemuk) \\ \midrule
0 & Rp 10.000.000 & Rp 10.000.000 \\
1 & Rp 11.200.000 & Rp 11.200.000 \\
2 & Rp 12.400.000 & Rp 12.544.000 \\
3 & Rp 13.600.000 & Rp 14.049.280 \\ \bottomrule
\end{tabular}
\end{center}
\textbf{Pertanyaan.} (a) Mengapa Bank B lebih besar meski bunganya sama?
(b) Apa pola pertumbuhannya? (c) Jika $10$ tahun, selisihnya membesar atau
mengecil?
\end{studikasus}

\subsection*{Jawaban lengkap studi kasus}

\begin{ingat}{Tunggal vs majemuk}
\textbf{Tunggal:} bunga selalu dihitung dari \emph{modal awal} $M_0$. Bunga tiap
tahun sama $\Rightarrow$ tumbuh \emph{linear} (barisan aritmetika).
\textbf{Majemuk:} bunga ditambahkan ke modal, lalu tahun berikut bunga dihitung
dari saldo yang sudah bertumbuh $\Rightarrow$ tumbuh \emph{eksponensial}
(barisan geometri).
\end{ingat}

\textbf{(a) Mengapa Bank B lebih besar?} Karena \emph{bunga berbunga}. Di Bank A,
tiap tahun bunga tetap $12\%\times10$jt $=1{,}2$jt (dari modal awal saja). Di
Bank B, bunga tahun ke-2 dihitung dari $11{,}2$jt (bukan $10$jt), sehingga
$12\%\times11{,}2$jt $=1{,}344$jt --- lebih besar. Selisih ini menumpuk tiap
tahun.

\textbf{(b) Pola pertumbuhan.} Bank A \emph{konstan}: naik tetap $+1{,}2$jt tiap
tahun (linear). Bank B \emph{makin cepat}: kenaikan $+1{,}2$; $+1{,}344$;
$+1{,}505$jt \dots\ (eksponensial).

\textbf{(c) Untuk $10$ tahun.} Selisih akan \textbf{membesar}. Pertumbuhan
eksponensial (Bank B) selalu akhirnya meninggalkan pertumbuhan linear (Bank A).
Bandingkan: Bank A $=10(1+10\cdot0{,}12)=22$jt; Bank B
$=10\cdot1{,}12^{10}\approx31{,}06$jt --- selisih $\approx9$jt, jauh lebih lebar
daripada $3$ tahun.

\begin{center}
\begin{tikzpicture}
\begin{axis}[width=11cm,height=5.5cm,xlabel=tahun $n$,ylabel=saldo (juta),
  xmin=0,xmax=10.5,ymin=9,ymax=33,legend pos=north west,legend style={font=\small},
  axis lines=left]
\addplot[very thick,biru,mark=*,mark size=1pt,domain=0:10,samples=11]{10*(1+0.12*x)};
\addlegendentry{Bank A (tunggal, linear)}
\addplot[very thick,oranye,mark=square*,mark size=1pt,domain=0:10,samples=11]{10*1.12^x};
\addlegendentry{Bank B (majemuk, eksponensial)}
\end{axis}
\end{tikzpicture}\\
{\small Garis lurus (tunggal) vs kurva melengkung naik (majemuk). Awalnya
berimpit, lalu menjauh --- ``keajaiban bunga majemuk''.}
\end{center}

% ############################################################
\section*{Subbab 9.1 — Bunga Tunggal}
\addcontentsline{toc}{section}{Subbab 9.1 — Bunga Tunggal}
% ############################################################

\begin{ingat}{Rumus bunga tunggal}
\[
M_n=M_0(1+n\,i),\qquad \text{bunga per periode}=M_0\,i,\qquad
\text{total bunga}=M_0\,n\,i.
\]
Tumbuh \emph{linear} (barisan aritmetika, beda $M_0 i$). Bunga selalu dari
$M_0$, bukan dari saldo yang bertumbuh.
\end{ingat}

% ---------------- SOAL 9.1.1 ----------------
\soalno{9.1.1}{nilai akhir bunga tunggal}
\begin{soalbox}{}
Rp$6.000.000$ ditabung dengan bunga tunggal $9\%$ per tahun selama $3$ tahun.
Hitung nilai akhirnya.
\end{soalbox}

\jawab
\langkah{1} Kenali: $M_0=6.000.000$, $i=0{,}09$, $n=3$.
\langkah{2} Rumus $M_n=M_0(1+ni)$:
\[
M_3=6.000.000\,(1+3\times0{,}09)=6.000.000\,(1+0{,}27)=6.000.000\times1{,}27.
\]
\langkah{3} $=7.620.000$.
\jadi nilai akhir $=$ Rp$7.620.000$ (total bunga $1.620.000=6.000.000\times3
\times0{,}09$).

\begin{jebakan}
\begin{itemize}[leftmargin=*,topsep=2pt]
  \item \textbf{Memakai rumus majemuk} $(1{,}09)^3$. Soal ini \emph{tunggal}:
  pakai $(1+ni)$, bukan pangkat.
  \item \textbf{Mengubah $9\%$ jadi $9$.} $i=0{,}09$ (bentuk desimal).
\end{itemize}
\end{jebakan}

\begin{variasi}
Ditanya total \emph{bunga} saja ($M_0 ni=1.620.000$); atau per bulan (konversi
$i$ dan $n$).
\end{variasi}

% ---------------- SOAL 9.1.2 ----------------
\soalno{9.1.2}{mencari modal awal}
\begin{soalbox}{}
Setelah $5$ tahun dengan bunga tunggal $6\%$ per tahun, nilai tabungan adalah
Rp$13.000.000$. Berapa modal awalnya?
\end{soalbox}

\jawab
\langkah{1} Diketahui $M_n=13.000.000$, $i=0{,}06$, $n=5$. Cari $M_0$.
\langkah{2} $M_n=M_0(1+ni)\Rightarrow 13.000.000=M_0(1+5\times0{,}06)
=M_0(1{,}3)$.
\langkah{3} $M_0=\dfrac{13.000.000}{1{,}3}=10.000.000$.
\jadi modal awal $=$ Rp$10.000.000$.

\begin{jebakan}
\begin{itemize}[leftmargin=*,topsep=2pt]
  \item \textbf{Mengurangi bunga secara kasar} ($13$jt $-30\%$). Bagi dengan
  faktor $(1+ni)=1{,}3$, jangan mengurangi persentase dari nilai akhir.
\end{itemize}
\end{jebakan}

\begin{variasi}
Mencari suku bunga $i$ (jika $M_0,M_n,n$ diketahui); mencari $M_0$ pada bunga
majemuk (bagi $(1+i)^n$).
\end{variasi}

% ---------------- SOAL 9.1.3 ----------------
\soalno{9.1.3}{mencari lama waktu}
\begin{soalbox}{}
Pada suku bunga tunggal $8\%$ per tahun, berapa lama hingga uang Rp$4.000.000$
menjadi Rp$5.280.000$?
\end{soalbox}

\jawab
\langkah{1} $M_n=M_0(1+ni)$: $5.280.000=4.000.000(1+0{,}08\,n)$.
\langkah{2} Bagi kedua ruas dengan $4.000.000$: $1{,}32=1+0{,}08n$.
\langkah{3} $0{,}32=0{,}08n\Rightarrow n=\dfrac{0{,}32}{0{,}08}=4$.
\jadi $n=4$ tahun.

\begin{jebakan}
\begin{itemize}[leftmargin=*,topsep=2pt]
  \item \textbf{Memakai logaritma.} Bunga \emph{tunggal} bersifat linear ---
  $n$ dicari dengan aljabar biasa, bukan logaritma (logaritma untuk majemuk).
\end{itemize}
\end{jebakan}

\begin{variasi}
Bunga majemuk (butuh logaritma, Soal 9.2.3); mencari $i$ dari $M_n,M_0,n$.
\end{variasi}

% ############################################################
\section*{Subbab 9.2 — Bunga Majemuk}
\addcontentsline{toc}{section}{Subbab 9.2 — Bunga Majemuk}
% ############################################################

\begin{ingat}{Rumus bunga majemuk}
\[
M_n=M_0(1+i)^n.
\]
Tumbuh \emph{eksponensial} (barisan geometri, rasio $1+i$). Untuk mencari
\emph{waktu}, gunakan logaritma: $n=\dfrac{\log(M_n/M_0)}{\log(1+i)}$.
\textbf{Konversi satuan:} bunga per bulan $=\dfrac{i_\text{tahun}}{12}$, dan $n$
dalam bulan. \textbf{Aturan 72:} waktu penggandaan $\approx\dfrac{72}{i\%}$.
\end{ingat}

% ---------------- SOAL 9.2.1 ----------------
\soalno{9.2.1}{nilai akhir bunga majemuk}
\begin{soalbox}{}
Rp$4.000.000$ ditabung dengan bunga majemuk $10\%$ per tahun selama $3$ tahun.
Hitung nilai akhirnya (gunakan $(1{,}1)^3=1{,}331$).
\end{soalbox}

\jawab
\langkah{1} $M_0=4.000.000$, $i=0{,}10$, $n=3$.
\langkah{2} $M_n=M_0(1+i)^n=4.000.000\,(1{,}1)^3=4.000.000\times1{,}331$.
\langkah{3} $=5.324.000$.
\jadi nilai akhir $=$ Rp$5.324.000$.
\emph{(Bandingkan tunggal: $4.000.000(1+3\cdot0{,}1)=5.200.000$. Selisih
$124.000$ dari efek bunga berbunga.)}

\begin{jebakan}
\begin{itemize}[leftmargin=*,topsep=2pt]
  \item \textbf{Memakai $(1+ni)=1{,}3$} (rumus tunggal). Majemuk memakai
  \emph{pangkat} $(1+i)^n=1{,}331$.
  \item \textbf{$4.000.000\times1{,}1\times3$.} Bukan dikali $3$; dipangkatkan
  $3$.
\end{itemize}
\end{jebakan}

\begin{variasi}
$n$ lebih besar memperbesar jurang tunggal--majemuk; mencari $M_0$ (bagi
$(1{,}1)^3$).
\end{variasi}

% ---------------- SOAL 9.2.2 ----------------
\soalno{9.2.2}{majemuk per bulan (konversi satuan)}
\begin{soalbox}{}
Modal Rp$2.000.000$ didepositokan dengan bunga majemuk $6\%$ per tahun
(diperhitungkan per bulan) selama $1$ tahun. Hitung nilai akhirnya.
\end{soalbox}

\begin{ingat}{Cocokkan satuan waktu}
``Per tahun diperhitungkan per bulan'' $\Rightarrow$ ubah bunga ke per bulan:
$i_\text{bulan}=\dfrac{6\%}{12}=0{,}5\%=0{,}005$, dan $n=12$ bulan.
\end{ingat}

\jawab
\langkah{1} $i_\text{bulan}=\dfrac{0{,}06}{12}=0{,}005$; $n=12$ (bulan).
\langkah{2} $M_{12}=2.000.000\,(1{,}005)^{12}$.
\langkah{3} $(1{,}005)^{12}\approx1{,}06168$, sehingga
$M_{12}\approx2.000.000\times1{,}06168\approx2.123.356$.
\jadi nilai akhir $\approx$ Rp$2.123.356$.
\emph{(Sedikit lebih besar daripada bunga tahunan biasa $2.000.000\times1{,}06
=2.120.000$, karena dimajemukkan lebih sering.)}

\begin{jebakan}
\begin{itemize}[leftmargin=*,topsep=2pt]
  \item \textbf{Memakai $i=0{,}06$ dengan $n=12$.} Itu mencampur satuan (bunga
  tahunan, periode bulanan). Bagi bunga dengan $12$ dulu.
  \item \textbf{Memakai $n=1$ dengan $i=0{,}005$.} Kalau bunga per bulan, $n$
  harus dalam bulan ($=12$).
\end{itemize}
\end{jebakan}

\begin{variasi}
Dimajemukkan per $3$ bulan (kuartal): $i=\frac{0{,}06}{4}$, $n=4$; bandingkan
frekuensi pemajemukan.
\end{variasi}

% ---------------- SOAL 9.2.3 ----------------
\soalno{9.2.3}{waktu penggandaan (logaritma \& Aturan 72)}
\begin{soalbox}{}
Pada bunga majemuk $8\%$ per tahun, berapa tahun hingga modal berlipat dua?
(Gunakan $\log 2=0{,}301$, $\log 1{,}08=0{,}0334$.)
\end{soalbox}

\begin{ingat}{Waktu = logaritma}
Berlipat dua berarti $M_n=2M_0$, yakni $(1+i)^n=2$. Karena $n$ ada di
\emph{pangkat}, gunakan logaritma (kaitan Bab 1):
$n=\dfrac{\log 2}{\log(1+i)}$.
\end{ingat}

\jawab
\langkah{1} Syarat berlipat dua: $(1{,}08)^n=2$.
\langkah{2} Ambil logaritma: $n\log1{,}08=\log2\Rightarrow
n=\dfrac{\log2}{\log1{,}08}=\dfrac{0{,}301}{0{,}0334}\approx9{,}01$.
\langkah{3} Jadi sekitar $9$ tahun.
\langkah{4} \textbf{Cek dengan Aturan 72:} $n\approx\dfrac{72}{i\%}=\dfrac{72}{8}
=9$ tahun --- cocok!
\jadi modal berlipat dua dalam $\approx9$ tahun.

\begin{jebakan}
\begin{itemize}[leftmargin=*,topsep=2pt]
  \item \textbf{Membagi $2$ dengan $0{,}08$} ($=25$). ``Berlipat dua'' memakai
  logaritma, bukan pembagian langsung --- $n$ terjebak di pangkat.
  \item \textbf{Salah bagi log.} $\frac{\log2}{\log1{,}08}$ (bukan $\log2-\log1{,}08$
  atau terbalik).
\end{itemize}
\end{jebakan}

\begin{variasi}
Berlipat \emph{tiga} ($(1{,}08)^n=3$, pakai $\log3$); Aturan 72 untuk suku bunga
lain ($9\%\Rightarrow8$ tahun).
\end{variasi}

% ============================================================
\begin{tips}
\textbf{Ringkasan strategi Bab 9.}
\begin{enumerate}[leftmargin=*,topsep=2pt]
  \item \emph{Kenali sistemnya:} \textbf{tunggal} $M_n=M_0(1+ni)$ (linear,
  aritmetika) vs \textbf{majemuk} $M_n=M_0(1+i)^n$ (eksponensial, geometri).
  \item \emph{Cocokkan satuan:} bunga tahunan $\to$ $n$ tahun; per bulan $\to$
  bagi $i$ dengan $12$ dan $n$ dalam bulan.
  \item \emph{Mencari waktu:} tunggal pakai aljabar; majemuk pakai
  \textbf{logaritma} $n=\frac{\log(M_n/M_0)}{\log(1+i)}$ (Aturan 72 untuk taksiran
  cepat penggandaan).
  \item \emph{Jangka panjang:} majemuk selalu menang jauh atas tunggal --- inilah
  jantung investasi (dan juga bahaya bunga pinjaman).
\end{enumerate}
\end{tips}

\end{document}
